\documentclass[11pt]{article}
\usepackage[letterpaper,margin=0.7in,headheight=15pt,headsep=18pt,footskip=25pt]{geometry}
\usepackage[T1]{fontenc}
\usepackage{helvet}
\renewcommand\familydefault{\sfdefault}
\usepackage{amsmath,amssymb,tikz,enumitem,fancyhdr,hyperref}
\hypersetup{colorlinks=true,urlcolor=blue,linkcolor=black,pdftitle={Week 65: Hyperbolic octagon streets - Adult guide},pdfauthor={Bellingham Math Circle}}
\setlength\parindent{0pt}\setlength\parskip{7pt}
\setlist[itemize]{leftmargin=17pt,itemsep=3pt,topsep=3pt}
\setlist[enumerate]{leftmargin=20pt,itemsep=4pt,topsep=3pt}
\pagestyle{fancy}\fancyhf{}\renewcommand\headrulewidth{0pt}
\fancyhead[L]{\small Week 65 / Hyperbolic octagon streets / Adult guide}
\fancyfoot[L]{\small Bellingham Math Circle / Week 65 / W65-FAC-v2}\fancyfoot[R]{\small\thepage}
\newcommand{\topic}[1]{\vspace{4pt}{\large\bfseries #1}\par}
\newcommand{\sub}[1]{{\bfseries #1}\par}
\newcommand{\route}[1]{\ensuremath{#1}}
\begin{document}
\topic{The mathematics first}
This is the regular hyperbolic tiling \(\{8,4\}\): each room has eight equal geodesic sides, each inside corner is a right angle, and four congruent rooms meet at every vertex. A geodesic is a straight route in this geometry. The Poincar\'e disk map preserves angles but distorts lengths. Its curved-looking streets are straight in the modeled world; its smaller-looking outer rooms are the same hyperbolic size as the central room.

\begin{enumerate}
\item \textbf{Local quarter-turns need not make a four-step loop.} In the ordinary square grid, moving one edge and turning left closes after four moves. Around one of these octagons, the same local rule first restores both position and heading after \textbf{eight} moves. After four the robot is at the opposite vertex. This concerns the supplied equal-edge walk, not every possible hyperbolic journey. Around two side-sharing octagons the outside walk has twelve quarter-turns and two straight junctions; two ordinary squares still need only four quarter-turns.
\item \textbf{The Euclidean unique-parallel rule fails.} The two complete straight-through streets drawn through \(P\) both miss the complete dashed street \(R\). This is a statement about entire lines, not just a finite crop. The exact check is on page~\pageref{complete-proof}. There are other hyperbolic lines through \(P\); the student task concerns the streets in this drawing.
\item \textbf{Routes and destinations are different things.} On the four-room map, two different two-door routes lead from Home to the star. There are eight four-door routes under the stated repeat-allowed rule. A square grid also has this shared-destination phenomenon. Here it prevents the mistaken idea that each outward route discovers a new room.
\end{enumerate}

\sub{What counts as evidence?}
Children experiment, predict, trace and explain with the maps. These experiments do not alone prove a statement about infinitely long lines. The adult checks use the circle construction and the actual room adjacency. All diagrams are original. Digital checks do not establish physical print usability or classroom success; this packet is \textbf{unpiloted}.

\topic{Readiness and a realistic first route}
\textbf{Grades 3--5, approximately.} Children need to follow a heading, distinguish a turn from moving, and check a partner's next branch. Reading can be shared. Counting to twelve is enough for the core; formal geometry, angles in degrees, multiplication and complex numbers are not prerequisites.

\textbf{First visit:} use Problems 1--2, then the two-room boundary comparison in Problem 3 if children still own the movement rule. Problems 4--5, the four-door census and adult growth extension can wait for another visit. Finishing every page is not the goal.

\textbf{Younger access:} a child may move an arrow and choose a local branch with an adult. If the adult supplies every turn, this is not an independent investigation. There is no separate K--1 packet; use a familiar tactile activity for that table rather than assuming this fills a whole K--5 session.
\newpage

\topic{Prepare the table}
\sub{Per pair, or a pair with a rotating referee}
\begin{itemize}
\item One student packet, six pages, printed single-sided on US Letter at 100\%. Offer one page at a time.
\item One slim arrow marker about 10 mm long, or an arrow drawn on a small scrap of paper; one second counter or coin; pencil and eraser. A second color helps overlay traces but is not required.
\item Some blank paper for route lists. The printed maps use line weights and labels and work without color. Keep the room-record sheet beside its board rather than flipping the board over.
\end{itemize}
For three working groups, prepare three packets plus a demonstration copy and one spare: \textbf{five packets, 30 sheets}. Have six working counters and two spares, three pencils plus a spare, and several blank sheets. No precision cutting, glue, compass, protractor, software or special model is needed.

Before children arrive, move a real marker on both octagon boards and check that it does not hide the outgoing branches or labels. On the two-room board, the closest junction centers are about 17.9 mm apart; a slim 10 mm arrow is preferable to a broad 15 mm disk. Rehearse the transition from street walking to crossing room sides. These physical checks have not yet been performed by the authors. If print is faint, darken it; do not shrink the maps to fit two pages on a sheet.

\topic{A short launch, then children take over}
Allow the usual five-minute run first. Take about three minutes for the shared demonstration using the visual at the top of student page 1.
\begin{enumerate}
\item Put an arrow at the example's starting dot. Say: ``One move takes us along one edge to the next dot.'' Move it, keeping its nose along the arriving edge.
\item Say: ``Now I can turn left, turn right, or keep going straight.'' Turn left at the arrival dot. The move has ended there; turning does not travel another edge.
\item Reset. Let a child make one move and choose a turn. Ask the partner to point to the next edge before the marker moves. Let both children try; stop the demonstration before completing the main loop.
\end{enumerate}
When introducing page 2, use its curved-edge example: keep the arrow along the road as it bends, then turn from the arrival direction. Following that printed bend does not count as a junction turn.

At a ready table, one child moves and the other checks. Swap after an attempt. In a group of three, rotate mover, checker and recorder/referee. An adult stays nearby but does not call out every branch. Page rotation and tracing with a finger are legitimate tools, not failures of mathematical ability.

\topic{A manageable hour}
After the launch, allow a few minutes of free legal moves. Spend roughly 15--20 minutes on Problems 1--2 and 10--15 minutes on the two-room comparison or another chosen investigation. Use the final five minutes to share one route or surprising observation. Save unfinished ideas for another meeting.

If attention fades, have a partner make a short legal route and the other replay it from the same position and heading. If keeping a heading remains too demanding, switch to the two-door room task after its own demonstration. Do not combine both movement rules at once.
\newpage

\topic{Problems 1--2: walking and turning}
\sub{Problem 1: ordinary square-grid routes}
A return includes the allowed turn at the final arrival at \(A\). The robot must then face along the original right-pointing arrow. No immediate reversal is allowed. The student asks for examples of three lengths, not all possible routes.

Here are checked examples. In this adult shorthand, \(E,N,W,S\) mean directions of \emph{traveled edges}; they are not turn commands:
\begin{center}
\begin{tabular}{cll}
Moves & Edge directions & Shape of the route \\
4 & \(E,N,W,S\) & one grid square \\
6 & \(E,E,N,W,W,S\) & a two-by-one rectangle \\
8 & \(E,E,N,N,W,W,S,S\) & a two-by-two square
\end{tabular}
\end{center}
All fit on the printed board. Each final southward arrival is followed by a left turn to face east. Other answers are possible, including a four-step loop traveled twice for eight moves. Repeated edges and early returns are not forbidden here.

\textbf{Hints, in order:}
\begin{enumerate}
\item ``Show me where the arrow will be after one move. Which turns are allowed there?''
\item ``Can your route go around a block and reach the starting dot? What way will the arrow face?''
\item ``What changes if the block is longer, or if you repeat your whole route?''
\end{enumerate}
Ask a ready pair what feature distinguishes a return to the dot from a return to the whole starting state. Do not require a written explanation after each attempt.

\sub{Problem 2: the octagon robot}
With left turns the traveled vertices are
\[A,B,C,D,E,F,G,H,A.\]
After four moves the robot is at \(E\), facing toward \(F\). The first full return is \textbf{eight moves}: at \(A\), the last left turn points toward \(B\). No earlier positive number of moves reaches \(A\).

With right turns the sequence is
\[A,H,G,F,E,D,C,B,A.\]
Again the fourth destination is \(E\), and the first full return is \textbf{eight}. Its final right turn points toward \(H\). Both answers count traveled edges, not the initial position.

\textbf{Hints, in order:} have the checker point to the arriving and departing edges at one corner; put a second counter at the start; replay slowly and stop after the fourth traveled edge. If ``left'' means the left side of the page to a child, turn the page until the arrow points away from them, then ask them to choose the branch.

The important comparison is a \emph{local} rule that children can enact. The octagon's corners really are right angles in this geometry. Its arcs bend on the printed map, so adding ordinary flat-paper heading changes as though every arc were a straight printed segment gives the wrong interpretation. Do not explain this as a faulty drawing or pretend it is an ordinary flat octagon.
\newpage

\topic{Problem 3: outside two rooms}
This is a purposeful continuation of the one-room walk. The octagon pair is viewed from the middle of its shared side, so both working rooms are readable. The new view preserves the geometry and local angles; it is not an ordinary Euclidean enlargement of one corner.

\sub{Answers on the two actual boards}
Start at the top endpoint of the shared wall and follow the indicated counterclockwise outside direction, keeping the two rooms on the left. The common wall is inside the cluster and is never part of this outside trip.
\begin{center}
\begin{tabular}{lrrr}
Board & Edge moves & Left quarter-turns & Straight junctions \\
Two ordinary squares & 6 & \textbf{4} & 2 \\
Two hyperbolic octagons & 14 & \textbf{12} & 2
\end{tabular}
\end{center}
The child task counts local turns, not two simultaneous tallies. Adults may note the traveled-edge totals afterwards. At the two endpoints of the shared wall the outside path goes straight: two right-angle room corners meet to make a straight angle. The midpoint of that wall is \emph{not} a junction and has no movement dot.

At the other twelve octagon boundary junctions there is one right-angle room corner, so the outside trip turns left. The final arrival at the top shared endpoint goes straight into the original heading. No last quarter-turn at the start is hidden in the answer.

\sub{A short explanation children can build}
One octagon has eight right-angle corners. Two octagons supply sixteen room-corner contributions. At each of the two ends of their shared wall, two of these corners join to make a straight passage. Those four contributions supply no outside quarter-turn. The twelve remaining corners each supply one: \(16-4=12\).

The same reasoning for two squares gives \(8-4=4\) quarter-turns. So the ordinary outside trip remains a rectangle with four turns, while the two-octagon trip needs twelve. For the optional edge count, the shared wall was counted once in each room but is not on the outside: \(8+8-2=14\) edges for the octagons and \(4+4-2=6\) for the squares.

\sub{Hints, in order}
\begin{enumerate}
\item ``Which wall is inside the pair? Can your finger stay on the outside all the way around?''
\item Stop at an endpoint of the shared wall. ``What is your arriving direction? Does the outside road turn here or keep going straight?''
\item ``Can you mark the places where you turn, so your partner can check the count? What happened to the two corners that meet at this end of the wall?''
\end{enumerate}
Do not count the apparent bending of a printed arc as a turn. A turn is a change to another local road direction at a marked junction. If counting interrupts the walk, mark the turning junctions first and count the marks afterwards; use this hint only when needed.

A successful comparison is more useful than a lecture about curvature. Let children keep the one-room and two-room records side by side and describe what changed. A general formula for larger clusters is not required.
\newpage

\topic{Problems 4--5: whole roads and room routes}
\sub{Problem 4: two streets through \(P\)}
There are \textbf{two} complete drawn streets through \(P\), and both miss the heavy dashed street \(R\). One bends across the upper-right part of the map; the other bends down its right side. Follow each in \emph{both} directions, taking the opposite branch at every crossing. A path that turns onto another street is a different object.

Ordinary Euclidean straight lines cannot have this entire crossing pattern: two distinct lines through one point cannot both be disjoint from a third line not through that point. On a finite piece of paper their intersections can be off the page, so ``I drew short segments that miss'' is not a counterexample. Complete lines, rather than short strokes, are the point of the comparison.

\textbf{Hints, in order:} cover distracting parts of the map and choose an incoming/outgoing pair at \(P\); trace that one street both ways; then trace the other in a different pencil style. Ask what would count as crossing \(R\). The picture supports noticing; page~\pageref{complete-proof} supplies the adult whole-line justification. Children need not reproduce that proof.

\sub{Problem 5: same destination, different journeys}
Demonstrate the separate room action with the printed \(C\to D\to E\) example. A move crosses one side into the neighboring room, never the central corner. Left/right street commands no longer apply. The four rooms are a recentered view of a corner of the same tiling, not four Euclidean squares.

The two two-door routes are \(H\to A\to *\) and \(H\to B\to *\), where \(H\) abbreviates Home in this adult key. The star is one shared room, not two copies. Repeated rooms and doors are allowed in the longer task. Its complete eight-route list is:
\begin{center}
\begin{tabular}{ll}
\(H\to A\to H\to A\to *\) & \(H\to B\to H\to A\to *\) \\
\(H\to A\to H\to B\to *\) & \(H\to B\to H\to B\to *\) \\
\(H\to A\to *\to A\to *\) & \(H\to B\to *\to A\to *\) \\
\(H\to A\to *\to B\to *\) & \(H\to B\to *\to B\to *\)
\end{tabular}
\end{center}
Why complete? The first room after Home is \(A\) or \(B\); the second is Home or the star; the third is \(A\) or \(B\); the fourth must be the star. These three independent two-way choices give all eight routes. An adult may record while children operate the coins.

\textbf{Hints, in order:} ask each partner to leave Home through a different side; replay the two routes and compare their final room; for the longer task, sort completed routes by their first visited room. Children can pool lists rather than each copying all eight. Stop after the two-door discovery if listing has replaced exploration. Ask what they would need to check before declaring a newly found route a newly found room.
\newpage

\topic{Why the model and complete-road claim are sound}\label{complete-proof}
This page is optional adult background, not a child prerequisite. Let the map be the open unit disk in the complex plane. Put
\[
 r=\sqrt{\sqrt2-1},\qquad
 v_k=r\exp\!\left(i\left(\frac{\pi}{8}+\frac{k\pi}{4}\right)\right),\quad k=0,\ldots,7.
\]
Join consecutive vertices by arcs of circles orthogonal to the unit boundary circle. When endpoints \(p,q\) are not on a common diameter, the supporting circle's center \(c\) satisfies
\[
 \operatorname{Re}(\overline c p)=\frac{|p|^2+1}{2},\qquad
 \operatorname{Re}(\overline c q)=\frac{|q|^2+1}{2};
 \qquad \rho^2=|c|^2-1.
\]
These equations determine the finite circle arcs in the builder. A geodesic through the disk center is a diameter: draw a straight segment and reflect in that line. For a finite supporting circle, side reflection is inversion in that circle. It produces congruent hyperbolic rooms, not ordinary Euclidean copies made by scaling the first room. At every vertex the incident circle tangents meet at right angles; four rooms fit smoothly around it.

\sub{A proof about the whole lines}
Let \(C=2^{1/4}\), and set \(P=v_0\). The supporting circles for the two streets through \(P\) have centers \(C e^{i\pi/4}\) and \(C\); the dashed street's circle has center \(-C\). All three have Euclidean radius \(r\).

Every point of the dashed circle has \(x\le -C+r<0\). Every point of the other two circles has, respectively,
\[
 x\ge C/\sqrt2-r>0,\qquad x\ge C-r>0.
\]
The first strict inequality follows from
\((C/\sqrt2)^2-r^2=1-\sqrt2/2>0\). Thus the entire dashed circle is separated from both other entire circles by the line \(x=0\). In particular, their arcs inside the disk cannot meet it. The two other arcs both pass through \(P\) and are distinct.

\sub{What ``straight'' means here}
At a crossing there are four local branches, separated by right angles. The opposite continuation lies on the same supporting circle or diameter and is the geodesic continuation. A person on this road makes no local turn. The printed arc's tangent direction changes along it because of the map. The boundary circle is an ideal boundary at infinite hyperbolic distance, not a street, a wall, or a collection of destinations.

\sub{Verification supplied with the source}
The student builder checks its exact diagram data and route lists. A separate standard-library checker reconstructs the tiling and deduplicates by both vertex sets and centers. It verifies eight sides per room, right-angle tangents, equal hyperbolic edge lengths, four rooms at every central corner, the whole-circle separation and layers \(1,8,48\). Numerical tolerances supplement the arguments above; they are not a proof of classroom usability.
\newpage

\topic{A return visit, and what to keep from it}
\sub{An optional growth question for a ready group}
The main packet does not ask children to count 48 tiny outer rooms. If children want to know how many rooms lie two door crossings from Home in the \emph{whole} tiling, start with the eight adjacent rooms. Each has seven doors other than the one back to Home. That produces \(8\times7=56\) non-backtracking two-door \emph{routes}.

At each of Home's eight corners, the two adjacent first-ring rooms share one second-ring room. Each such destination is counted twice. There are eight distinct duplicated destinations, and the checked local geometry has no further overlaps. Hence there are \(56-8=48\) distinct rooms exactly two crossings away: 40 reached once and eight reached twice. Let children reason from one enlarged corner and symmetry, rather than label the entire tiny second ring.

An ordinary square-room grid has exact-distance layers \(1,4,8\); this octagon tiling starts \(1,8,48\). These values are a concrete comparison of early growth. They do \textbf{not} justify multiplying by seven forever or prove an asymptotic growth rate. Later overlaps matter.

Other useful continuations are to replay a robot route with its handedness reversed, compare two views of the same marked junction, or organize the eight four-door routes so a partner can check completeness. Keep a question alive when children still have mathematical choices; a second visit is preferable to rushing three representations.

\sub{After the meeting}
Record which problems children actually tried, not just which pages were offered. Note whether they could choose local branches without the adult, which diagrams needed enlargement, and whether they distinguished walking \emph{along} a street from crossing \emph{between} rooms. Keep one route, sketch or quoted observation. Label a claim as tried, conjectured, checked in some cases, or explained. Use those observations to revise the next session.

\topic{Sources and scope}
\begin{itemize}
\item David E. Joyce, \emph{Hyperbolic Tessellations}, Clark University. Tiling notation and the Poincar\'e map. \href{https://mathcs.clarku.edu/~djoyce/poincare/poincare.html}{mathcs.clarku.edu/\textasciitilde djoyce/poincare/poincare.html}
\item Kathryn Mann, \emph{DIY Hyperbolic Geometry}, Mathcamp 2015, Days 1--2. A precedent for doing and drawing before formalism; its paper model is an approximation. \href{https://e.math.cornell.edu/people/mann/papers/DIYhyp.pdf}{e.math.cornell.edu/people/mann/papers/DIYhyp.pdf}
\item Cornell Mathematics Explorers' Club, \emph{Geometries of Surfaces}. The octagon surface example there uses \(\{8,8\}\), with 45-degree corners. \href{https://pi.math.cornell.edu/~mec/Winter2009/Victor/part4.htm}{pi.math.cornell.edu/\textasciitilde mec/Winter2009/Victor/part4.htm}
\end{itemize}
Sources checked October 5, 2026. Wording, diagrams and the particular packet are newly authored around established mathematics. This \(\{8,4\}\) tiling takes inspiration from an uncertain octagon-seminar recollection; it does not claim to recover that seminar. Week 64's Euclidean edge-identified octagon is a separate activity. No external figures or passages are reproduced.
\end{document}
